\chapter{Comparación de funciones}\label{ch:b2:comparison}

El análisis asintótico ---el arte de sustituir una cantidad
complicada por otra sencilla más un error controlado--- se inició en
el volumen del primer año con los desarrollos de Taylor. Este capítulo
lo convierte en disciplina propia: desarrollos sobre escalas
generales, la comparación serie--integral con toda su fuerza
asintótica, la fórmula de Stirling (demostrada por completo) y el
estudio sistemático de las sucesiones definidas implícitamente. Estas
técnicas son el pan de cada día del análisis asintótico, y todo
capítulo posterior que estime algo ---series, integrales,
probabilidades--- come de esta mesa.

\section{Relaciones de comparación y escalas}

\begin{definition}\label{def:b2:comparison:landau}
Cerca de un punto $a$ ($a \in \R$ o $\pm\infty$), para funciones (o
para sucesiones, con $n \to \infty$): $f = o(g)$, $f = O(g)$, $f
\sim g$ como en el volumen del primer año. Una \emph{escala de
comparación}\index{escala de comparación} en $a$ es una familia de
funciones positivas, comparables dos a dos y totalmente ordenada por
$o(\cdot)$; la escala estándar en $+\infty$ es
\[
x^{\alpha} (\ln x)^{\beta}
\qquad (\alpha, \beta \in \R),
\]
ordenada lexicográficamente en $(\alpha, \beta)$ y refinada cuando
hace falta con exponenciales $\eu^{\gamma x}$.
\end{definition}

\begin{definition}[Desarrollo asintótico]\label{def:b2:comparison:expansion}
$f$ admite el \emph{desarrollo asintótico}\index{desarrollo
asintótico}
\[
f = c_1 \varphi_1 + c_2\varphi_2 + \dots + c_k \varphi_k +
o(\varphi_k)
\qquad (\varphi_{i+1} = o(\varphi_i) \text{ en la escala})
\]
cuando los restos sucesivos cumplen las estimaciones indicadas.
Entonces los coeficientes son únicos: $c_1 = \lim f/\varphi_1$ y,
por inducción, $c_{i+1} = \lim\,(f - \sum_{j \leq i}
c_j\varphi_j)/\varphi_{i+1}$.
\end{definition}

\begin{example}
Los desarrollos de Taylor son \omterm{def:b2:comparison:expansion}{desarrollos asintóticos} sobre la
escala $(x - a)^k$ en $a$. Pero la noción es estrictamente más
amplia: en $+\infty$,
\[
\frac{1}{x - \ln x}
= \frac1x \cdot \frac{1}{1 - \frac{\ln x}{x}}
= \frac1x + \frac{\ln x}{x^2} + o\Bigl(\frac{\ln x}{x^2}\Bigr),
\]
un desarrollo sobre la escala mixta: no se aplica ningún teorema de
Taylor, solo el desarrollo geométrico y el cálculo con $o$.
\end{example}

\begin{example}[La escala estándar está realmente ordenada]\label{ex:b2:comparison:scaleorder}
La afirmación lexicográfica de la~\cref{def:b2:comparison:landau} necesita una línea de
demostración por caso. Comparemos $x^{\alpha}(\ln x)^{\beta}$ y
$x^{\alpha'}(\ln x)^{\beta'}$ en $+\infty$. Si $\alpha <
\alpha'$: el cociente es $x^{\alpha - \alpha'}(\ln x)^{\beta -
\beta'} \to 0$, porque una potencia negativa de $x$ aplasta
cualquier potencia de $\ln x$ (póngase $x = \eu^t$:
$\eu^{(\alpha - \alpha')t}\,t^{\beta - \beta'} \to 0$ por el
límite exponencial-gana-a-polinomio del volumen del primer año).
Si $\alpha = \alpha'$ y $\beta < \beta'$: el cociente es $(\ln
x)^{\beta - \beta'} \to 0$ directamente. Así pues, los pares
$(\alpha, \beta)$, ordenados lexicográficamente, ordenan la
escala por $o(\cdot)$; y la sustitución $x = \eu^t$ es el truco
universal para las comparaciones mixtas de potencias y
logaritmos.
\end{example}

\begin{example}[Ordenar una fauna variada]\label{ex:b2:comparison:ranking}
Las escalas han de estar \emph{ordenadas}; he aquí el ejercicio
estándar. En $+\infty$, comparemos $n^{10}$,
$\eu^{\sqrt{\ln n}\,\cdot\,\sqrt n}$, $2^n$ y $n^{\ln n}$ tomando
logaritmos:
\[
10\ln n
\;\ll\; (\ln n)^2
\;\ll\; \sqrt{n\ln n}
\;\ll\; n\ln 2 ,
\]
donde $a_n \ll b_n$ significa $a_n = o(b_n)$; la segunda entrada es
$\ln(n^{\ln n})$. Las exponenciales conservan esos saltos estrictos
(si $\ln u_n - \ln v_n \to -\infty$, entonces $u_n/v_n \to 0$), de
modo que
\[
n^{10} = o\bigl(n^{\ln n}\bigr),
\qquad
n^{\ln n} = o\bigl(\eu^{\sqrt{n\ln n}}\bigr),
\qquad
\eu^{\sqrt{n\ln n}} = o(2^n) .
\]
La moraleja, por partida doble: compárese siempre a través de
logaritmos (diferencias de logaritmos, no cocientes de logaritmos),
y no se concluya nunca $u_n \sim v_n$ a partir de $\ln u_n \sim \ln
v_n$; el par $n^{10}$ y $n^{\ln n}$ tiene cociente de logaritmos
que tiende a $\infty$, pero $2^n$ y $4^n$ tienen cociente de
logaritmos exactamente $2$ y son tremendamente no equivalentes.
\end{example}

\section{Comparación serie--integral, en clave asintótica}

\begin{theorem}\label{thm:b2:comparison:seriesintegral}
Sea $f$ \omterm{def:b2:metric:continuity}{continua}, positiva y \emph{decreciente} sobre
$\intco{1}{+\infty}$.
\begin{enumerate}
  \item Si $\int_1^{\infty} f$ converge, los restos cumplen
        \[
        \int_{n+1}^{\infty} f \;\leq\; \sum_{k > n} f(k) \;\leq\;
        \int_{n}^{\infty} f .
        \]
  \item Si $\int_1^\infty f$ diverge, las sumas parciales cumplen
        $\sum_{k=1}^{n} f(k) = \int_1^n f + C + o(1)$ para cierta
        constante $C$: la diferencia $\sum_{k \leq n} f(k) - \int_1^n
        f$ \emph{converge}.
\end{enumerate}
\end{theorem}

\begin{proof}
El encuadre $f(k+1) \leq \int_k^{k+1} f \leq f(k)$ (por ser
decreciente) fue el recurso del primer año; sumando sobre $k \geq
n+1$, o bien sobre $k \geq n$, se obtiene (1). Para (2), pongamos
$u_k = f(k) - \int_k^{k+1} f$: por el encuadre, $0 \leq u_k \leq
f(k) - f(k+1)$, de modo que las sumas parciales de $\sum u_k$ están
acotadas por la suma telescópica $f(1) - f(n+1) \leq f(1)$: la serie
converge. Además, la sucesión $\bigl(\int_n^{n+1} f\bigr)_n$ es no
creciente ($f$ decrece) y no negativa, luego convergente.
Escribiendo
\[
\sum_{k=1}^{n} f(k) - \int_1^n f
= \sum_{k=1}^{n} u_k + \int_n^{n+1} f ,
\]
el miembro derecho converge cuando $n \to \infty$: la diferencia
converge a una constante $C$, que es el enunciado (2).
\end{proof}

\begin{example}[El desarrollo armónico]\label{ex:b2:comparison:harmonic}
Para $f(t) = \frac1t$: $H_n = \ln n + \gamma + o(1)$, recuperando
la constante de Euler (volumen del primer año) con una demostración
más limpia. Un orden más allá (\cref{exo:b2:comparison:3}):
\[
H_n = \ln n + \gamma + \frac{1}{2n} + o\Bigl(\frac1n\Bigr).
\]
Los números hacen visible la ganancia en $n = 10$: $H_{10} =
2.928968\dots$ y $\ln 10 = 2.302585\dots$, de modo que la
estimación bruta de $\gamma$ es $H_{10} - \ln 10 = 0.626383$, con
un error de $0.049$; restando la corrección $\frac1{20}$ se obtiene
$0.576383$, que dista de $\gamma = 0.577216$ solo
$8.3\cdot10^{-4}$, cantidad que es a su vez el término siguiente
$\frac{1}{12\cdot100}$ del desarrollo, como demuestra el problema
de fin de semana (pregunta 8).
\end{example}

\begin{example}[Un $\ln(n!)$ tosco, sin Stirling]\label{ex:b2:comparison:lnfactorial}
El recurso del encuadre basta ya para localizar $\ln(n!)$. Como
$\ln$ es creciente,
\[
\int_{k-1}^{k}\ln t\,\dd t \;\leq\; \ln k \;\leq\;
\int_{k}^{k+1}\ln t\,\dd t ,
\]
y sumando sobre $k = 2, \dots, n$ (con $\int_1^n\ln = n\ln n - n +
1$):
\[
n\ln n - n + 1 \;\leq\; \ln(n!) \;\leq\; (n+1)\ln(n+1) - n .
\]
Ambas vallas valen $n\ln n - n + O(\ln n)$, de donde $\ln(n!) =
n\ln n - n + O(\ln n)$ y, en particular, $\ln(n!) \sim n\ln n$. Lo
que Stirling añade son los dos peldaños siguientes ---el
$\frac12\ln n$ y la constante $\ln\sqrt{2\pi}$---, que cuestan la
suma telescópica más fina del~\cref{thm:b2:comparison:stirling}. Saber qué precisión compra cada
herramienta es la mitad del oficio asintótico.
\end{example}

\begin{example}[Las cancelaciones exigen desarrollos]\label{ex:b2:comparison:cancellation}
Calculemos el límite de $\sqrt{n^2 + n} - n$. Ambos términos son
$\sim n$, y ``$\sim n - n$'' no significa nada: los equivalentes no
se pueden restar. Desarrollemos en su lugar:
\[
\sqrt{n^2 + n} - n
= n\Bigl(\sqrt{1 + \tfrac1n} - 1\Bigr)
= n\Bigl(\frac{1}{2n} - \frac{1}{8n^2} +
O\Bigl(\frac1{n^3}\Bigr)\Bigr)
= \frac12 - \frac{1}{8n} + O\Bigl(\frac{1}{n^2}\Bigr) :
\]
el límite es $\frac12$, con la velocidad de aproximación
$\frac1{8n}$ de regalo. El mecanismo merece nombre: una diferencia
de dos cantidades grandes equivalentes vive por completo en sus
términos \emph{siguientes}, así que hay que desarrollar hasta el
primer orden en el que ambos miembros difieren, arrastrando el
resto para certificar que a ese orden no sobrevive nada más.
\end{example}

\begin{example}[Una comparación divergente, resuelta]\label{ex:b2:comparison:loglog}
Para $f(t) = \frac{1}{t\ln t}$ sobre $\intco{2}{+\infty}$
(\omterm{def:b2:metric:continuity}{continua}, positiva y decreciente): $\int_2^x f = \ln\ln x -
\ln\ln 2 \to \infty$, luego por el~\cref{thm:b2:comparison:seriesintegral}~(2),
\[
\sum_{k=2}^{n}\frac{1}{k\ln k} = \ln\ln n + C + o(1)
\]
para cierta constante $C$. Dos lecciones. Primera: la divergencia
es real pero glacial; la suma parcial supera $4$ por primera vez
hacia $n \approx \eu^{\eu^{4 - C}}$, astronómicamente grande.
Segunda: la \emph{forma} $\ln\ln n$ la ha entregado una primitiva,
no una conjetura; para términos monótonos, la integral es el
dispositivo canónico de suma, y la constante $C$ ---como la
$\gamma$ de Euler--- es la memoria de los primeros términos.
\end{example}

\section{La fórmula de Stirling}

\begin{lemma}[Integrales de Wallis, revisitadas]\label{lem:b2:comparison:wallis}
Sea $W_n = \int_0^{\pi/2} \sin^n t\,\dd t$. Entonces $nW_nW_{n-1} =
\frac\pi2$ para $n \geq 1$, la sucesión $(W_n)$ decrece y $W_n \sim
\sqrt{\dfrac{\pi}{2n}}$.
\end{lemma}

\begin{proof}
La integración por partes da $nW_n = (n-1)W_{n-2}$ ($n \geq 2$),
luego $nW_nW_{n-1}$ es constante en $n$ e igual a $1 \cdot W_1 W_0 =
\frac\pi2$. Decrecimiento: $\sin^{n+1} \leq \sin^n$ sobre
$\intcc{0}{\frac\pi2}$. El encaje, en detalle: la monotonía da
$W_{n+1} \leq W_n \leq W_{n-1}$ y, dividiendo por $W_{n-1} > 0$,
\[
\frac{n}{n+1} = \frac{W_{n+1}}{W_{n-1}} \leq
\frac{W_n}{W_{n-1}} \leq 1 ,
\]
donde la identidad de la izquierda sale de la recurrencia en el
índice $n + 1$. Ambas cotas tienden a $1$: $W_n \sim W_{n-1}$, de
donde
\[
nW_n^2 \sim nW_nW_{n-1} = \frac\pi2
\qquad\Longrightarrow\qquad
W_n \sim \sqrt{\frac{\pi}{2n}} .
\]
\end{proof}

\begin{example}[Las primeras integrales de Wallis]\label{ex:b2:comparison:wallisvalues}
A partir de $W_0 = \frac\pi2$, $W_1 = 1$ y la recurrencia $nW_n =
(n-1)W_{n-2}$:
\[
W_2 = \frac\pi4, \qquad
W_3 = \frac23, \qquad
W_4 = \frac{3\pi}{16}, \qquad
W_5 = \frac{8}{15}, \qquad
W_6 = \frac{5\pi}{32}.
\]
Los índices pares llevan un $\pi$ y los impares son racionales: los
dos productos entrelazados de las formas cerradas. Numéricamente,
$W_6 \approx 0.4909$ frente al valor asintótico $\sqrt{\pi/12}
\approx 0.5116$: ya en $n = 6$ el equivalente está a menos del
$5\,\%$, y la identidad del producto es exacta para todo $n$:
$6\,W_6W_5 = 6\cdot\frac{5\pi}{32}\cdot\frac8{15} = \frac\pi2$.
Tablas pequeñas como esta son la manera más barata de cazar un
desliz algebraico antes de que infecte un argumento asintótico.
\end{example}

\begin{theorem}[Stirling]\label{thm:b2:comparison:stirling}
\[
n! \;\sim\; \sqrt{2\pi n}\, \Bigl(\frac{n}{\eu}\Bigr)^{\!n}
.\index{fórmula de Stirling}
\]
\end{theorem}

\begin{proof}
\emph{Paso 1: $n! \sim C \sqrt n\, (n/\eu)^n$ para cierta constante
$C > 0$.} Pongamos
\[
d_n = \ln(n!) - \Bigl(n + \frac12\Bigr)\ln n + n .
\]
Entonces
\[
d_n - d_{n+1}
= \Bigl(n + \frac12\Bigr) \ln\frac{n+1}{n} - 1
= \Bigl(n + \frac12\Bigr)\Bigl(\frac1n - \frac{1}{2n^2} +
\frac{1}{3n^3} + o\bigl(n^{-3}\bigr)\Bigr) - 1
= \frac{1}{12n^2} + o\Bigl(\frac{1}{n^2}\Bigr),
\]
por el desarrollo de Taylor de $\ln(1 + \frac1n)$. La serie $\sum
(d_n - d_{n+1})$ converge, pues, absolutamente (comparación con
$\sum n^{-2}$), de modo que $(d_n)$ converge, digamos a $d$;
exponenciando, $n! \sim C\sqrt n\,(n/\eu)^n$ con $C = \eu^{d}$.

\emph{Paso 2: $C = \sqrt{2\pi}$ mediante Wallis.} La forma cerrada
$W_{2p} = \frac{(2p)!}{4^p (p!)^2}\cdot\frac\pi2$ (procedente de la
recurrencia, cálculo del primer año rehecho en el marco del~\cref{lem:b2:comparison:wallis}) se combina con el paso 1:
\[
W_{2p} \sim \frac{C\sqrt{2p}\,(2p/\eu)^{2p}}
{4^p\,\bigl(C\sqrt p\,(p/\eu)^p\bigr)^2}\cdot\frac{\pi}{2}
= \frac{\sqrt{2p}}{C\,p}\cdot\frac{\pi}{2}
= \frac{\pi}{C}\cdot\frac{1}{\sqrt{2p}} .
\]
Comparando con $W_{2p} \sim \sqrt{\frac{\pi}{4p}}$
(\cref{lem:b2:comparison:wallis}): la igualdad
$\frac{\pi}{C\sqrt{2p}} = \sqrt{\frac{\pi}{4p}}\,(1 + o(1))$ obliga
a $C = \pi \sqrt{\frac{4p}{2p\,\pi}} = \sqrt{2\pi}$.
\end{proof}

\begin{example}[Coeficiente binomial central]\label{ex:b2:comparison:centralbinomial}
\[
\binom{2n}{n} = \frac{(2n)!}{(n!)^2}
\sim \frac{\sqrt{4\pi n}\,(2n/\eu)^{2n}}{2\pi n\,(n/\eu)^{2n}}
= \frac{4^n}{\sqrt{\pi n}} :
\]
la probabilidad de que un paseo aleatorio simétrico vuelva a $0$ en
el instante $2n$ es $\sim \frac{1}{\sqrt{\pi n}}$, un anuncio del~\cref{ch:b2:randomvar}.
\end{example}

\begin{remark}[Perspectivas dentro de este volumen]\label{rem:b2:comparison:perspectives}
Todos los capítulos cuantitativos que vienen hablan la lengua de
este. El~\cref{ch:b2:series} clasifica las series comparando sus
términos con la escala $n^{-\alpha}(\ln n)^{-\beta}$, y su problema
de fin de semana cartografía por completo esa frontera. El~\cref{ch:b2:integration} hace lo mismo con las integrales
impropias, con la escala idéntica en la variable \omterm{def:b2:metric:continuity}{continua}. El~\cref{ch:b2:powerseries} calcula radios de convergencia a partir de
$\limsup\abs{a_n}^{1/n}$, un ejercicio de equivalentes de raíces
$n$-ésimas donde Stirling es la clave habitual ($\sqrt[n]{n!} \sim
\frac n\eu$, \cref{exo:b2:comparison:4}). Y los capítulos de
probabilidad cobran Stirling directamente: las estimaciones locales
del~\cref{ch:b2:randomvar} para coeficientes binomiales son el~\cref{ex:b2:comparison:centralbinomial} y el~\cref{ex:b2:comparison:lopsided} palabra por palabra. La asintótica
no es aquí un capítulo: es el acento del volumen.
\end{remark}

\begin{method}[Lista de comprobación del arranque sucesivo]\label{met:b2:comparison:checklist}
Antes de fiarse de un desarrollo obtenido por arranques sucesivos,
audítense cuatro puntos. (1) \emph{Primero la existencia:} la raíz
o la sucesión ha de quedar fijada (monotonía, valores intermedios)
antes de todo desarrollo; los símbolos sin referente se desarrollan
maravillosamente y no significan nada. (2) \emph{Un orden por
pasada:} cada sustitución solo es fiable hasta el orden de la
estimación introducida; extraer dos términos nuevos de una sola
pasada es la fuente clásica de coeficientes erróneos. (3)
\emph{Los restos viajan con nosotros:} arrástrese el $o(\cdot)$ por
todos los pasos algebraicos y déjese la absorción (los términos
menores tragados por los restos mayores) para el final, de forma
explícita. (4) \emph{Auditoría numérica:} evalúese en un valor
honesto de $n$; un error de coeficiente sobrevive a una nueva
deducción algebraica con sorprendente frecuencia, y casi nunca
sobrevive a la aritmética.
\end{method}

\begin{remark}[Errores frecuentes]\label{rem:b2:comparison:pitfalls}
(i) Los equivalentes se \emph{suman} mal: de $u_n \sim n + \ln n$ y
$v_n \sim -n$ \emph{no} puede concluirse $u_n + v_n \sim \ln n$;
las cancelaciones exigen desarrollos con restos explícitos y nunca
equivalentes a secas. (ii) No hay que exponenciar nunca una
equivalencia: $n + 1 \sim n$ pero $\eu^{n+1} \not\sim \eu^n$; la
dirección segura es tomar logaritmos de equivalentes que tienden a
$+\infty$ (problema de fin de semana de este capítulo, pregunta
24). (iii) Un \omterm{def:b2:comparison:expansion}{desarrollo asintótico} va unido a una \emph{escala}:
escribir $f = \frac1x + o\bigl(\frac1{x^2}\bigr)$ afirma más que $f
= \frac1x + o\bigl(\frac1x\bigr)$, y mezclar ambas cosas invalida
el álgebra posterior. (iv) En los arranques sucesivos hay que
sustituir el desarrollo actual \emph{entero}, resto incluido;
descartar un $o(\cdot)$ a mitad de pasada produce coeficientes
plausibles pero falsos. (v) La comparación serie--integral necesita
la monotonía: para términos oscilantes falla sin más (compárese
$\sum\frac{\sin k}k$, \cref{ch:b2:series}).
\end{remark}

\begin{example}[Stirling en números]\label{ex:b2:comparison:stirlingnum}
Para $n = 10$: la fórmula da $\sqrt{20\pi}\,(10/\eu)^{10} \approx
3\,598\,696$ frente a $10! = 3\,628\,800$, con error relativo
$8.3\cdot10^{-3}$, notable para un enunciado ``asintótico'' en $n =
10$. El error tiene estructura ---el refinamiento exacto $n! =
\sqrt{2\pi n}\,(n/\eu)^n\bigl(1 + \frac1{12n} + O(n^{-2})\bigr)$---
y su primera corrección $\frac1{120} \approx 8.3\cdot10^{-3}$
explica casi exactamente la discrepancia observada. La maquinaria
de Euler--Maclaurin del problema de fin de semana es precisamente
la fuente sistemática de esos términos correctores.
\end{example}

\begin{remark}[Dónde se usa este capítulo]
La comparación asintótica es la gramática de todo lo cuantitativo
que viene después: los criterios de convergencia y el panorama de
Bertrand del~\cref{ch:b2:series}, los criterios de integrabilidad
del~\cref{ch:b2:integration}, los cálculos de radio de convergencia
del~\cref{ch:b2:powerseries} y los teoremas límite del~\cref{ch:b2:randomvar} (donde Stirling mueve las estimaciones de de
Moivre--Laplace). El volumen del tercer año industrializa la única
idea que aquí demostramos a mano ---extraer el término principal y
acotar el resto--- convirtiéndola en el método de Laplace y en la
convergencia dominada.
\end{remark}

\begin{example}[Una integral comparada consigo misma: $\int_2^x \frac{\dd t}{\ln t}$]\label{ex:b2:comparison:liworked}
La caja de herramientas de comparación también funciona con
integrales. Sea $F(x) = \int_2^x\frac{\dd t}{\ln t}$ (el integrando
es \omterm{def:b2:metric:continuity}{continuo} sobre $\intco2\infty$). Integrando por partes:
\[
F(x) = \Bigl[\frac{t}{\ln t}\Bigr]_2^x +
\int_2^x\frac{\dd t}{(\ln t)^2}
= \frac{x}{\ln x} + O\Bigl(\int_2^x\frac{\dd t}{(\ln
t)^2}\Bigr) + O(1),
\]
y la integral restante es $o\bigl(\frac{x}{\ln x}\bigr)$:
pártase en $\sqrt x$ y acótese mediante
\[
\int_2^{\sqrt x}\frac{\dd t}{(\ln t)^2} \leq \sqrt x
\qquad\text{y}\qquad
\int_{\sqrt x}^{x}\frac{\dd t}{(\ln t)^2} \leq
\frac{x}{(\ln\sqrt x)^2} = \frac{4x}{(\ln x)^2} .
\]
Por tanto $F(x) \sim \frac{x}{\ln x}$. Quien haya visto el teorema
de los números primos en el problema de fin de semana de este
capítulo reconocerá a $F$: es la integral logarítmica, el mejor
estimador de $\pi(x)$, y el cálculo muestra que coincide con
$\frac{x}{\ln x}$ en primer orden.
\end{example}

\begin{example}[Stirling sobre un binomial desequilibrado]\label{ex:b2:comparison:lopsided}
La misma rutina de tres factoriales que en el~\cref{ex:b2:comparison:centralbinomial} da, para $\binom{3n}{n} =
\frac{(3n)!}{n!\,(2n)!}$:
\[
\binom{3n}{n} \sim
\frac{\sqrt{6\pi n}\,(3n/\eu)^{3n}}
{\sqrt{2\pi n}\,(n/\eu)^{n}\cdot\sqrt{4\pi n}\,(2n/\eu)^{2n}}
= \sqrt{\frac{3}{4\pi n}}\,
\Bigl(\frac{27}{4}\Bigr)^{\!n} .
\]
La tasa exponencial $\frac{27}4 = \frac{3^3}{2^2}$ es $\eu^{3n\,
H(1/3)}$ en la notación de entropía de la teoría de la información:
los binomiales desequilibrados crecen estrictamente más despacio
que el central $4^n$ por cada dos pasos ---aquí, $(27/4)^{1/3}
\approx 1.89 < 2$ por paso---. Toda asintótica de binomiales en
combinatoria y en probabilidad (\cref{ch:b2:randomvar}) es este
mismo cálculo con otros pesos.
\end{example}

\section{Sucesiones definidas implícitamente}

\begin{method}\label{met:b2:comparison:implicit}
Para hallar la asintótica de las soluciones $x_n$ de una ecuación
$F(x, n) = 0$:\index{sucesión implícita}
\begin{enumerate}
  \item \emph{Localizar}: demuéstrese la existencia y la unicidad de
        $x_n$ en un intervalo determinado (monotonía, teorema del
        valor intermedio) y hállese su comportamiento tosco (límite,
        orden de crecimiento).
  \item \emph{Arrancar}: sustitúyase la forma tosca $x_n =
        (\text{término principal})(1 + \varepsilon_n)$ en la
        ecuación y despéjese el orden siguiente de $\varepsilon_n$;
        repítase, refinando un orden por pasada.
\end{enumerate}
\end{method}

\begin{example}\label{ex:b2:comparison:tan}
Para $n \geq 1$, la ecuación $\tan x = x$ tiene exactamente una
solución $x_n$ en $\intoo{n\pi - \frac\pi2}{n\pi + \frac\pi2}$ (la
función $\tan x - x$ crece allí de $-\infty$ a $+\infty$, pues su
derivada es $\tan^2 x \geq 0$). \emph{Tosco:} $x_n = n\pi +
\frac\pi2 - y_n$ con $y_n \in \intoo{0}{\pi}$; como $x_n \to \infty$
y $\tan x_n = x_n \to +\infty$, $x_n$ se acerca a la asíntota por la
izquierda: $y_n \to 0$. \emph{Arranque:} $\tan x_n = \cot y_n =
\frac{1}{\tan y_n} \sim \frac{1}{y_n}$, y la ecuación $\cot y_n =
x_n \sim n\pi$ da $y_n \sim \frac{1}{n\pi}$. Por tanto
\[
x_n = n\pi + \frac\pi2 - \frac{1}{n\pi} +
o\Bigl(\frac1n\Bigr),
\]
y el proceso continúa hasta cualquier orden
(\cref{exo:b2:comparison:6}).
\end{example}

\begin{example}[Una segunda vuelta del método]\label{ex:b2:comparison:xplusln}
Resolvamos asintóticamente $x + \ln x = n$. \emph{Localizar:} $x
\mapsto x + \ln x$ crece de $-\infty$ a $+\infty$ sobre
$\intoo{0}{+\infty}$: hay una única raíz $x_n$, y $x_n \to \infty$.
\emph{Tosco:} de $\ln x_n = o(x_n)$ resulta $x_n \sim n$.
\emph{Arranque:} de $x_n = n - \ln x_n$ y $\ln x_n = \ln n + o(1)$
(logaritmos de equivalentes, con ambos miembros $\to \infty$):
\[
x_n = n - \ln n + o(1) ;
\]
una pasada más, con $\ln x_n = \ln\bigl(n - \ln n + o(1)\bigr) =
\ln n - \frac{\ln n}{n} + o\bigl(\frac{\ln n}n\bigr)$:
\[
x_n = n - \ln n + \frac{\ln n}{n} +
o\Bigl(\frac{\ln n}{n}\Bigr).
\]
(Comprobación en $n = 100$: la raíz es $x \approx 95.4415$; la
fórmula de tres términos da $100 - 4.6052 + 0.0461 = 95.4409$ y la
de dos términos $95.3948$; cada pasada gana el orden previsto.)
Mismo bucle, tercer paisaje: al método del~\cref{met:b2:comparison:implicit} le da igual qué aspecto tenga la
ecuación, con tal de que cada pasada aísle la incógnita dominante.
\end{example}

\section{Ejercicios}

\begin{exercise}[$\star$]\label{exo:b2:comparison:1}
Desarrolla en $+\infty$, dos términos más allá del principal:
\[
\sqrt{x^2 + x + 1} ,
\qquad
\ln(x^2 + x) - 2\ln x,
\qquad
\frac{x + \sin x}{x - \ln x} .
\]
\end{exercise}

\begin{exercise}[$\star$]\label{exo:b2:comparison:2}
Da el carácter (convergencia o divergencia) y, en caso divergente,
la asintótica principal de $\sum_{k \leq n} k^\alpha$ para $\alpha >
-1$, $\alpha = -1$ y $\alpha < -1$, mediante el~\cref{thm:b2:comparison:seriesintegral}.
\end{exercise}

\begin{exercise}[$\star\star$]\label{exo:b2:comparison:3}
Demuestra que $H_n = \ln n + \gamma + \frac{1}{2n} +
o\bigl(\frac1n\bigr)$. \emph{(Estudia $v_n = H_n - \ln n - \gamma$:
prueba que $v_n - v_{n+1} = \frac{1}{2n^2} + O(n^{-3})$ y suma la
cola, comparando con $\sum_{k \geq n} \frac{1}{2k^2} \sim
\frac{1}{2n}$, \cref{thm:b2:comparison:seriesintegral}~(1).)}
\end{exercise}

\begin{exercise}[$\star\star$]\label{exo:b2:comparison:4}
Usando Stirling, halla equivalentes de $\dfrac{(3n)!}{(n!)^3}$;
$\;\dfrac{n!}{n^n}$; $\;\sqrt[n]{n!}$ (en la forma $\frac
n\eu(1 + o(1))$, precisada a dos términos).
\end{exercise}

\begin{exercise}[$\star\star$]\label{exo:b2:comparison:5}
Para $n \geq 2$, demuestra que $x^n + x = 1$ tiene una única
solución $x_n \in \intoo{0}{1}$, que $x_n \to 1$, y establece
\[
x_n = 1 - \frac{\ln n}{n} + o\Bigl(\frac{\ln n}{n}\Bigr).
\]
\emph{(A partir de $x_n^n = 1 - x_n$: toma logaritmos y arranca con
$x_n = 1 - \varepsilon_n$.)}
\end{exercise}

\begin{exercise}[$\star\star$]\label{exo:b2:comparison:6}
Lleva el~\cref{ex:b2:comparison:tan} un orden más allá:
\[
x_n = n\pi + \frac\pi2 - \frac{1}{n\pi} + \frac{1}{2n^2\pi} +
o\Bigl(\frac{1}{n^2}\Bigr).
\]
\emph{(Escribe $\cot y_n = x_n$ exactamente, desarrolla $\cot y =
\frac1y - \frac y3 + o(y)$ y $x_n = n\pi(1 + \frac{1}{2n} -
\dots)$, e identifica.)}
\end{exercise}

\begin{exercise}[$\star\star$]\label{exo:b2:comparison:7}
Determina $\lim_{n\to\infty} \dfrac{1}{n!}\sum_{k=0}^{n} k!$
\emph{(acota la suma de todos los términos salvo los dos últimos)} y
deduce el \omterm{def:b2:comparison:expansion}{desarrollo asintótico} $\sum_{k \leq n} k! = n!\bigl(1 +
\frac1n + O(n^{-2})\bigr)$.
\end{exercise}

\begin{exercise}[$\star\star\star$]\label{exo:b2:comparison:8}
Sea $u_0 > 0$ y $u_{n+1} = u_n + \dfrac{1}{u_n}$. Demuestra que $u_n
\to \infty$, después que $u_n \sim \sqrt{2n}$ \emph{(estudia
$u_n^2$: sus incrementos son $2 + u_n^{-2}$; suma)}, y refina:
\[
u_n = \sqrt{2n}\Bigl(1 + \frac{\ln n}{8n} +
o\Bigl(\frac{\ln n}{n}\Bigr)\Bigr).
\]
\emph{(De $u_n^2 = 2n + \sum_{k<n} u_k^{-2} + u_0^2$ y $u_k^2 \sim
2k$: la suma vale $\sim \frac12\ln n$ por el~\cref{thm:b2:comparison:seriesintegral}.)}
\end{exercise}

\begin{exercise}[$\star\star\star$]\label{exo:b2:comparison:9}
(Una suma de Riemann con truco) Determina el comportamiento
asintótico de
\[
S_n = \sum_{k=1}^{n} \frac{1}{n + k\ln n} .
\]
\emph{(Saca factor $n$: $S_n = \frac1n\sum_k \bigl(1 + \frac{k\ln
n}{n}\bigr)^{-1}$; reconoce una suma de tipo Riemann con un
parámetro $t = \ln n$ que varía lentamente, calcula $\int_0^1
\frac{\dd u}{1 + tu} = \frac{\ln(1+t)}{t}$ y concluye $S_n \sim
\frac{\ln\ln n}{\ln n}$.)}
\end{exercise}

\begin{exercise}[$\star$]\label{exo:b2:comparison:10}
Demuestra la identidad $(\ln n)^{\ln n} = n^{\ln\ln n}$ y ordena
después las siguientes expresiones en orden creciente por $o(\cdot)$
en el infinito, con demostración: $n^2$, $(\ln n)^{\ln n}$, $2^n$,
$n!$, $n^n$.
\end{exercise}

\begin{exercise}[$\star\star$]\label{exo:b2:comparison:11}
(Cola de $\sum 1/k^2$, dos términos) Usando la suma telescópica
exacta $\sum_{k > n} \frac{1}{k(k+1)} = \frac{1}{n+1}$ y la
descomposición $\frac1{k^2} = \frac{1}{k(k+1)} +
\frac{1}{k^2(k+1)}$, demuestra que
\[
\sum_{k > n} \frac{1}{k^2}
= \frac1n - \frac{1}{2n^2} + O\Bigl(\frac{1}{n^3}\Bigr).
\]
\end{exercise}

\begin{exercise}[$\star\star\star$]\label{exo:b2:comparison:12}
Sea $u_0 = \frac12$ y $u_{n+1} = u_n + \eu^{-u_n}$. Demuestra que
$u_n \to \infty$ y, poniendo $v_n = \eu^{u_n}$ y probando que
$v_{n+1} = v_n + 1 + \frac{1}{2v_n} + O\bigl(v_n^{-2}\bigr)$,
establece
\[
u_n = \ln n + \frac{\ln n}{2n} + O\Bigl(\frac1n\Bigr).
\]
\end{exercise}

\section{Problema: arranques sucesivos, de Euler--Maclaurin a los
primos}

Una cantidad implícita o acumulada rara vez entrega su asintótica de
una vez; se extrae en pasadas, cada una realimentando la estimación
previa en la relación que la define. Este problema de fin de semana
entrena ese bucle sobre ecuaciones nuevas, demuestra la \emph{fórmula
de Euler--Maclaurin} de primer orden (la mejora trapezoidal de la
comparación serie--integral, con barras de error rigurosas), invierte
$x\ln x = n$ y cobra el cheque más famoso del método: a partir del
teorema de los números primos, admitido, la ley asintótica $p_n \sim
n\ln n$ del $n$-ésimo primo.

\begin{problem}[{Problema de fin de semana --- la corrección de
Euler--Maclaurin y la asintótica del $n$-ésimo primo}]
\label{pb:b2:comparison:1}

\medskip
\noindent\textbf{Parte I --- El bucle de arranque sobre una
ecuación nueva.}
\begin{enumerate}
  \item Demuestra la afirmación de unicidad de la~\cref{def:b2:comparison:expansion}: si $f = \sum_{i\leq k}
        c_i\varphi_i + o(\varphi_k) = \sum_{i \leq k}
        c_i'\varphi_i + o(\varphi_k)$ sobre la misma escala,
        entonces $c_i = c_i'$ para todo $i$. Lleva después el
        ejemplo mixto del curso un peldaño más allá:
        \[
        \frac{1}{x - \ln x} = \frac1x + \frac{\ln x}{x^2} +
        \frac{(\ln x)^2}{x^3} + o\Bigl(\frac{(\ln
        x)^2}{x^3}\Bigr) \qquad (x \to +\infty),
        \]
        y explica por qué no aparece ningún término
        $\frac{c}{x^2}$.
  \item Prueba que para todo $n \geq 1$ la ecuación $\eu^x + x = n$
        tiene exactamente una solución real $x_n$, y que $x_n \to
        +\infty$ con $x_n \sim \ln n$.
  \item Arranca dos veces:
        \[
        x_n = \ln n - \frac{\ln n}{n} - \frac{(\ln n)^2}{2n^2} +
        o\Bigl(\frac{(\ln n)^2}{n^2}\Bigr).
        \]
  \item Comprueba numéricamente en $n = 1000$: compara $x_{1000}
        \approx 6.90083$ con los valores de uno, dos y tres términos
        de la pregunta 3, con cinco decimales.
\end{enumerate}

\medskip
\noindent\textbf{Parte II --- Euler--Maclaurin, orden uno.}
\begin{enumerate}[resume]
  \item Demuestra la identidad del núcleo trapezoidal: para $g$ de
        clase $C^2$ sobre $\intcc{0}{1}$,
        \[
        \int_0^1 g(t)\,\dd t = \frac{g(0) + g(1)}{2}
        - \frac12\int_0^1 t(1 - t)\,g''(t)\,\dd t
        \]
        \emph{(integra $\frac12 t(1-t)g''$ por partes dos veces)}.
  \item Sea $f$ de clase $C^2$ sobre $\intco{1}{+\infty}$ con
        $\int_1^\infty \abs{f''} < \infty$. Prueba que
        \[
        E_n = \sum_{k=1}^{n} f(k) - \int_1^n f -
        \frac{f(1) + f(n)}{2}
        \]
        converge a una constante $E$, con la cota de la cola
        $\abs{E - E_n} \leq \frac18\int_n^\infty\abs{f''}$: la
        \emph{fórmula de Euler--Maclaurin} de primer orden.
  \item Aplícalo a $f(t) = \frac1t$: demuestra que
        \[
        H_n = \ln n + \gamma + \frac{1}{2n} + \varepsilon_n,
        \qquad \abs{\varepsilon_n} \leq \frac{1}{8n^2},
        \]
        reforzando el~\cref{exo:b2:comparison:3} (identifica la
        constante con $\gamma$ comparando con el~\cref{ex:b2:comparison:harmonic}).
  \item Extrae el coeficiente siguiente: prueba que $\varepsilon_n
        = -\frac{1}{12n^2} + o\bigl(\frac1{n^2}\bigr)$ \emph{(los
        incrementos de $E_n$ son $\frac12\int_0^1t(1-t)f''(n+t)\dd t
        = \frac1{12}f''(n) + o(f''(n))$; suma la cola con el~\cref{thm:b2:comparison:seriesintegral})}.
  \item Aplica la pregunta 6 a $f = \ln$: vuelve a deducir en tres
        líneas la convergencia de $d_n = \ln n! - (n + \frac12)\ln n
        + n$ (paso 1 del~\cref{thm:b2:comparison:stirling}), con la
        velocidad de error de regalo $d_n = d + O\bigl(\frac1n\bigr)$.
  \item Aplica la pregunta 6 a $f(t) = \frac{1}{\sqrt t}$: prueba
        que
        \[
        \sum_{k=1}^{n}\frac1{\sqrt k} = 2\sqrt n + c +
        \frac{1}{2\sqrt n} + O\Bigl(\frac{1}{n^{3/2}}\Bigr)
        \]
        para cierta constante $c$, y evalúa todos los términos en $n
        = 10^4$ (la constante es $c \approx -1.4604$).
\end{enumerate}

\medskip
\noindent\textbf{Parte III --- Inversión: la ecuación $x\ln x =
n$.}
\begin{enumerate}[resume]
  \item Prueba que $x\ln x = n$ tiene exactamente una solución $x_n
        \in \intco{1}{+\infty}$ para $n \geq 1$, que $x_n \to
        \infty$ y que $\ln x_n \sim \ln n$.
  \item Deduce la inversión de un término $x_n \sim
        \dfrac{n}{\ln n}$ y arranca una vez más:
        \[
        \ln x_n = \ln n - \ln\ln n + o(1),
        \qquad
        x_n = \frac{n}{\ln n}\Bigl(1 + \frac{\ln\ln n}{\ln n}
        + o\Bigl(\frac{\ln\ln n}{\ln n}\Bigr)\Bigr).
        \]
  \item Comprueba en $n = 10^6$: la raíz verdadera es $x \approx
        87\,848$; compárala con los valores de un término ($\approx
        72\,382$) y de dos términos ($\approx 86\,140$), y explica
        la lentitud de la ganancia (el parámetro del desarrollo es
        $\frac{\ln\ln n}{\ln n}$, que solo vale $\approx 0.19$ en $n
        = 10^6$).
  \item \emph{Admitamos} ahora el teorema de los números primos: el
        número $\pi(x)$ de primos $\leq x$ cumple $\pi(x) \sim
        \frac{x}{\ln x}$ cuando $x \to \infty$ (demostrado
        honestamente en el volumen del tercer año). Escribiendo
        $p_n$ para el $n$-ésimo primo, justifica $\pi(p_n) = n$ y
        aplica la inversión de las preguntas 11--12 para demostrar
        que
        \[
        p_n \sim n \ln n .
        \]
  \item Dividendos: (a) prueba que $\sum_{k \leq n} p_k \sim
        \frac{n^2\ln n}{2}$ \emph{(compara $\sum k\ln k$ con $\int
        t\ln t\,\dd t$)}; (b) calcula la probabilidad aproximada de
        que un entero elegido al azar uniformemente con $100$ cifras
        sea primo ($\ln 10^{100} \approx 230.26$: alrededor de uno
        de cada $230$).
\end{enumerate}

\medskip
\noindent\textbf{Parte IV --- El método exportado: $x\tan x = 1$.}
\begin{enumerate}[resume]
  \item Prueba que para cada $n \geq 1$ la ecuación $\tan x =
        \frac1x$ tiene exactamente una solución $x_n$ en
        $\intoo{n\pi}{\,n\pi + \frac\pi2}$, y que $z_n = x_n - n\pi
        \to 0^+$.
  \item Un término: $z_n \sim \dfrac{1}{n\pi}$.
  \item Prueba que el desarrollo de $z_n$ \emph{no} tiene término
        $\frac{c}{n^2}$: $z_n = \frac1{n\pi} +
        O\bigl(\frac{1}{n^3}\bigr)$.
  \item Tres términos: usando $\arctan u = u - \frac{u^3}3 +
        O(u^5)$ y $\frac1{x_n} = \frac{1}{n\pi} -
        \frac{z_n}{(n\pi)^2} + O(n^{-3}\cdot z_n^2)$, demuestra que
        \[
        x_n = n\pi + \frac{1}{n\pi} -
        \frac{4}{3\pi^3 n^3} + o\Bigl(\frac{1}{n^3}\Bigr).
        \]
  \item Comprueba en $n = 3$: la raíz verdadera es $x_3 \approx
        9.5293344$; compara los valores de uno y de tres términos, y
        contrasta en una frase con el $\tan x = x$ del curso
        (\cref{ex:b2:comparison:tan}): dónde se sitúa cada sucesión
        dentro de su ventana, y por qué.
\end{enumerate}

\medskip
\noindent\textbf{Parte V --- Un arranque dinámico, reglas del
juego, síntesis.}
\begin{enumerate}[resume]
  \item Sea $u_0 \in \intoo{0}{\pi}$ y $u_{n+1} = \sin u_n$. Prueba
        que $u_n \to 0$ de manera decreciente y calcula el límite de
        $\dfrac{1}{u_{n+1}^2} - \dfrac{1}{u_n^2}$ \emph{(desarrolla
        $\sin^{-2}$ mediante $\sin u = u - \frac{u^3}6 + o(u^3)$)}.
  \item Deduce, mediante las medias de Cesàro (volumen del primer
        año), el clásico
        \[
        u_n \sim \sqrt{\frac{3}{n}} .
        \]
  \item (Numérica certificada) Usando la cota rigurosa de la
        pregunta 7, prueba que evaluar $\ln n + \gamma + \frac1{2n}$
        en $n = 10^6$ da $H_{10^6}$ con un error de a lo sumo
        $1.25\cdot10^{-13}$: una suma de un millón de términos
        calculada con trece cifras mediante tres términos.
  \item (Reglas del juego) Demuestra o refuta, con demostraciones o
        contraejemplos: (a) si $u_n \sim v_n \to +\infty$, entonces
        $\ln u_n \sim \ln v_n$; (b) si $u_n \sim v_n$, entonces
        $\eu^{u_n} \sim \eu^{v_n}$; (c) si $f \sim g$ en $+\infty$
        ($f, g$ derivables), entonces $f' \sim g'$.
  \item (Síntesis) Una frase para cada punto: el bucle de arranque
        del~\cref{met:b2:comparison:implicit} tal como se usa en las
        partes I, III y IV; qué añade la corrección trapezoidal al~\cref{thm:b2:comparison:seriesintegral}; por qué la
        inversión de $x\ln x$ es exactamente el puente de $\pi(x)$ a
        $p_n$; y cuál de las reglas de la pregunta 24 protegió cada
        paso. Nombra las dos cumbres: la fórmula de Euler--Maclaurin
        (de primer orden) y la ley asintótica del $n$-ésimo primo.
\end{enumerate}
\end{problem}
